\documentclass{article}
\usepackage{amsmath,amssymb,amsthm}
\usepackage{algorithm}
\usepackage[noend]{algpseudocode}
\usepackage{bm}
\usepackage{hyperref}
\newtheorem{theorem}{Theorem}
\newtheorem{lemma}{Lemma}
\newtheorem{assumption}{Assumption}
\newcommand{\EE}{\mathbb{E}}
\newcommand{\RR}{\mathbb{R}}
\newcommand{\norm}[1]{\left\|#1\right\|}
\newcommand{\inner}[1]{\left\langle #1 \right\rangle}
\begin{document}

\section*{Quantized Stochastic Gradient Descent (QSGD)}

\section*{Mathematical Formulation}
Let $f:\RR^{d}\to\RR$ be a convex, $L$‑smooth objective.  
At iteration $t$, a stochastic gradient $g_t$ is drawn:
\[
g_t \;=\; \nabla f(w_t;\xi_t),\qquad 
\EE_{\xi_t}[\,g_t\,|\,w_t]=\nabla F(w_t),\qquad 
\EE_{\xi_t}\!\left[\|g_t-\nabla F(w_t)\|^{2}\mid w_t\right]\le\sigma^{2},
\]
where $F(w)=\EE_{\xi}[f(w;\xi)]$ is the population loss.

A (possibly biased) quantization operator $Q:\RR^{d}\to\RR^{d}$ satisfies
\[
\EE_{\mathcal{Q}}[\,Q(g)\mid g]=g,\qquad 
\EE_{\mathcal{Q}}\!\left[\|Q(g)-g\|^{2}\mid g\right]\le \omega\|g\|^{2},
\tag{1}
\]
with quantization variance factor $\omega\ge0$.

The update rule of **Quantized SGD** is
\[
\boxed{ \; w_{t+1}= w_t - \eta_t\, Q(g_t)\; } \tag{2}
\]
where $\{\eta_t\}_{t\ge0}$ is a (deterministic) stepsize schedule.

\section*{Pseudocode}
\begin{algorithm}[H]
\caption{Quantized Stochastic Gradient Descent}
\begin{algorithmic}[1]
\Require Initial point $w_{0}\in\RR^{d}$, stepsizes $\{\eta_t\}$, quantizer $Q(\cdot)$
\For{$t=0,1,\dots,T-1$}
    \State Sample a mini‑batch $\xi_t$ and compute stochastic gradient $g_t=\nabla f(w_t;\xi_t)$
    \State Quantize: $\tilde g_t = Q(g_t)$
    \State Update: $w_{t+1}= w_t - \eta_t \tilde g_t$
\EndFor
\State \Return $ \displaystyle \bar w_T=\frac1T\sum_{t=0}^{T-1} w_t$
\end{algorithmic}
\end{algorithm}

\section*{Dry Run Example}
Consider the scalar quadratic
\[
f(w)=(w-3)^{2},\qquad \nabla f(w)=2(w-3).
\]
We use a deterministic “quantizer’’ $Q(g)=g$ (i.e. $\omega=0$) for illustration,
stepsize $\eta=0.1$, and start from $w_{0}=0$.

\begin{center}
\begin{tabular}{c|c|c|c}
$t$ & $g_t=2(w_t-3)$ & $w_{t+1}=w_t-\eta g_t$ & $f(w_{t+1})$ \\ \hline
0 & $2(0-3)=-6$ & $0-0.1(-6)=0.6$ & $(0.6-3)^2=5.76$\\
1 & $2(0.6-3)=-4.8$ & $0.6-0.1(-4.8)=1.08$ & $(1.08-3)^2=3.6864$\\
2 & $2(1.08-3)=-3.84$ & $1.08-0.1(-3.84)=1.464$ & $(1.464-3)^2=2.3650$\\
\end{tabular}
\end{center}
The objective value decreases at each iteration, illustrating the algorithmic flow.

\newpage
\section*{Assumptions}
\begin{assumption}[Convexity and Smoothness]\label{ass:convex}
The population loss $F$ is convex and $L$‑smooth:
\[
F(w')\ge F(w)+\inner{\nabla F(w)}{w'-w},\qquad 
\|\nabla F(w)-\nabla F(w')\|\le L\|w-w'\|,\;\forall w,w'.
\]
\end{assumption}
\begin{assumption}[Stochastic Gradient Bounded Variance]\label{ass:sgdvar}
For all $t$,
\[
\EE\!\left[\|g_t-\nabla F(w_t)\|^{2}\mid w_t\right]\le\sigma^{2}.
\]
\end{assumption}
\begin{assumption}[Unbiased Quantization with Bounded Relative Variance]\label{ass:quant}
The quantizer $Q$ satisfies (1) with a constant $\omega\ge0$.
\end{assumption}
\begin{assumption}[Stepsize]\label{ass:stepsize}
The stepsizes are non‑increasing and satisfy
\[
\eta_t\le \frac{1}{L(1+\omega)}\quad\forall t .
\]
\end{assumption}

\section*{Main Convergence Theorem}
\begin{theorem}[Expected Suboptimality of QSGD]\label{thm:main}
Under Assumptions~\ref{ass:convex}–\ref{ass:stepsize}, for any $T\ge1$,
\[
\frac1T\sum_{t=0}^{T-1}\EE\!\big[\,F(w_t)-F(w^{\star})\,\big]
\;\le\;
\frac{\norm{w_0-w^{\star}}^{2}}{2\eta_T T}
\;+\;
\frac{\eta_{\max}}{2}\,( \sigma^{2}+ \omega L^{2} D^{2}),
\tag{3}
\]
where $w^{\star}\in\arg\min F$, $D:=\sup_{t}\norm{w_t-w^{\star}}$, and $\eta_{\max}:=\max_{t}\eta_t$.
In particular, with a constant stepsize $\eta_t=\eta\le\frac{1}{L(1+\omega)}$,
\[
\EE\!\big[F(\bar w_T)-F(w^{\star})\big]\le
\frac{\norm{w_0-w^{\star}}^{2}}{2\eta T}
\;+\;
\frac{\eta}{2}\big(\sigma^{2}+\omega L^{2}D^{2}\big).
\tag{4}
\]
Choosing $\eta=\Theta\!\big(\frac{1}{\sqrt{T}}\big)$ yields the optimal rate
\[
\EE\!\big[F(\bar w_T)-F(w^{\star})\big]=\mathcal{O}\!\Big(\frac{1}{\sqrt{T}}\Big).
\]
\end{theorem}

\section*{Theoretical Properties}
\begin{itemize}
\item \textbf{Stability.} The condition $\eta\le 1/(L(1+\omega))$ guarantees that the quantization noise does not destabilise the recursion (see Step~[5] of the proof).
\item \textbf{Hyper‑parameter Sensitivity.} The bound (3) shows a linear dependence on $\eta$ for the variance term and an inverse dependence for the bias term; thus the usual $\eta\propto 1/\sqrt{T}$ trade‑off remains valid even with quantization.
\item \textbf{Comparison to vanilla SGD.} Setting $\omega=0$ recovers the classical SGD bound (cf. Example 1). Hence quantization incurs an additional additive term $\frac{\eta\omega L^{2}D^{2}}{2}$.
\item \textbf{Special Cases.}
\begin{enumerate}
\item \emph{Deterministic gradients} ($\sigma=0$) – only the quantization variance $\omega$ remains.
\item \emph{Exact communication} ($\omega=0$) – the bound collapses to the standard SGD result.
\end{enumerate}
\end{itemize}

\section*{Discussion}
The theorem confirms that unbiased quantization with bounded relative variance preserves the $\mathcal{O}(1/\sqrt{T})$ convergence rate of stochastic gradient methods. The extra term scales with $\omega$ and the smoothness constant $L$, reflecting the intuition that finer quantization (smaller $\omega$) yields behaviour closer to unquantized SGD. Practically, one can select $\omega$ by adjusting the number of bits; the stepsize condition (Assumption~\ref{ass:stepsize}) is the only modification required in the learning‑rate schedule.

\newpage
\section*{Preliminaries (Definitions and Lemmas)}
\begin{lemma}[Three‑Point Identity]\label{lem:three}
For any vectors $a,b,c\in\RR^{d}$ and any $\eta>0$,
\[
\|a-c\|^{2}= \|a-b\|^{2}+2\inner{b-c}{a-b}+\|b-c\|^{2}.
\]
\end{lemma}
\begin{lemma}[Smoothness Upper Bound]\label{lem:smooth}
If $F$ is $L$‑smooth then for all $x,y$,
\[
F(y)\le F(x)+\inner{\nabla F(x)}{y-x}+\frac{L}{2}\|y-x\|^{2}.
\]
\end{lemma}

\section*{Main Proof (Step‑by‑Step Derivation)}
We prove Theorem~\ref{thm:main}.  Throughout, expectations are taken w.r.t. all randomness up to time $t$ unless otherwise noted.

\noindent\textbf{Step 1.} Start from the squared distance recursion using (2) and Lemma~\ref{lem:three} with $a=w_{t+1}$, $b=w_t$, $c=w^{\star}$:
\[
\|w_{t+1}-w^{\star}\|^{2}
= \|w_t-w^{\star}\|^{2}
-2\eta_t\inner{Q(g_t)}{w_t-w^{\star}}
+\eta_t^{2}\|Q(g_t)\|^{2}. \tag{5}
\]

\noindent\textbf{Step 2.} Take conditional expectation $\EE_t[\cdot]\equiv\EE[\cdot\mid w_t]$.
Using unbiasedness of $Q$ (Assumption~\ref{ass:quant}) we have
\[
\EE_t[Q(g_t)]=\EE_t[g_t]=\nabla F(w_t). \tag{6}
\]
Hence
\[
\EE_t\!\left[\inner{Q(g_t)}{w_t-w^{\star}}\right]
= \inner{\nabla F(w_t)}{w_t-w^{\star}}. \tag{7}
\]

\noindent\textbf{Step 3.} Bound the second moment of $Q(g_t)$.  
From (1) and the law of total expectation,
\[
\EE_t\!\big[\|Q(g_t)\|^{2}\big]
= \EE_t\!\big[\|g_t\|^{2}\big] + \EE_t\!\big[\|Q(g_t)-g_t\|^{2}\big]
\le \EE_t\!\big[\|g_t\|^{2}\big] + \omega\,\EE_t\!\big[\|g_t\|^{2}\big]
= (1+\omega)\EE_t\!\big[\|g_t\|^{2}\big]. \tag{8}
\]

\noindent\textbf{Step 4.} Decompose $\EE_t[\|g_t\|^{2}]$ into bias and variance components.
\[
\EE_t\!\big[\|g_t\|^{2}\big]
= \|\nabla F(w_t)\|^{2} + \EE_t\!\big[\|g_t-\nabla F(w_t)\|^{2}\big]
\le \|\nabla F(w_t)\|^{2}+\sigma^{2}. \tag{9}
\]

\noindent\textbf{Step 5.} Substitute (7)–(9) into (5) and take full expectation.
\[
\EE\!\big[\|w_{t+1}-w^{\star}\|^{2}\big]
\le \EE\!\big[\|w_t-w^{\star}\|^{2}\big]
-2\eta_t\EE\!\big[\inner{\nabla F(w_t)}{w_t-w^{\star}}\big]
+\eta_t^{2}(1+\omega)\EE\!\big[\|\nabla F(w_t)\|^{2}+\sigma^{2}\big].
\tag{10}
\]

\noindent\textbf{Step 6.} Relate the inner product to function suboptimality using convexity (Assumption~\ref{ass:convex}):
\[
\inner{\nabla F(w_t)}{w_t-w^{\star}} \ge F(w_t)-F(w^{\star}). \tag{11}
\]

\noindent\textbf{Step 7.} Apply smoothness to bound $\|\nabla F(w_t)\|^{2}$.  
From $L$‑smoothness and convexity,
\[
\|\nabla F(w_t)\|^{2}\le 2L\big(F(w_t)-F(w^{\star})\big). \tag{12}
\]
(Proof: use Lemma~\ref{lem:smooth} with $y=w^{\star}$ and rearrange.)

\noindent\textbf{Step 8.} Plug (11) and (12) into (10):
\[
\EE\!\big[\|w_{t+1}-w^{\star}\|^{2}\big]
\le \EE\!\big[\|w_t-w^{\star}\|^{2}\big]
-2\eta_t\EE\!\big[F(w_t)-F(w^{\star})\big]
+2\eta_t^{2}L(1+\omega)\EE\!\big[F(w_t)-F(w^{\star})\big]
+\eta_t^{2}(1+\omega)\sigma^{2}. \tag{13}
\]

\noindent\textbf{Step 9.} Use the stepsize restriction (Assumption~\ref{ass:stepsize}) i.e. $\eta_t\le \frac{1}{L(1+\omega)}$,
which guarantees $2\eta_t^{2}L(1+\omega)\le \eta_t$.
Hence the coefficient of $\EE[F(w_t)-F(w^{\star})]$ in (13) becomes
\[
-2\eta_t + 2\eta_t^{2}L(1+\omega)\;\le\; -\eta_t . \tag{14}
\]

\noindent\textbf{Step 10.} Inequality (13) simplifies to
\[
\EE\!\big[\|w_{t+1}-w^{\star}\|^{2}\big]
\le \EE\!\big[\|w_t-w^{\star}\|^{2}\big]
-\eta_t\EE\!\big[F(w_t)-F(w^{\star})\big]
+\eta_t^{2}(1+\omega)\sigma^{2}. \tag{15}
\]

\noindent\textbf{Step 11.} Rearrange (15) to isolate the suboptimality term:
\[
\eta_t\EE\!\big[F(w_t)-F(w^{\star})\big]
\le \EE\!\big[\|w_t-w^{\star}\|^{2}\big]
- \EE\!\big[\|w_{t+1}-w^{\star}\|^{2}\big]
+\eta_t^{2}(1+\omega)\sigma^{2}. \tag{16}
\]

\noindent\textbf{Step 12.} Sum (16) over $t=0,\dots,T-1$ and divide by $T$:
\[
\frac1T\sum_{t=0}^{T-1}\eta_t\EE\!\big[F(w_t)-F(w^{\star})\big]
\le \frac{\norm{w_0-w^{\star}}^{2}}{T}
+\frac{1}{T}\sum_{t=0}^{T-1}\eta_t^{2}(1+\omega)\sigma^{2}. \tag{17}
\]

\noindent\textbf{Step 13.} Define $\eta_{\max}:=\max_{t}\eta_t$ and use $\eta_t\le\eta_{\max}$:
\[
\frac1T\sum_{t=0}^{T-1}\EE\!\big[F(w_t)-F(w^{\star})\big]
\le \frac{\norm{w_0-w^{\star}}^{2}}{2\eta_T T}
+\frac{\eta_{\max}}{2}( \sigma^{2}+ \omega L^{2}D^{2}),
\tag{18}
\]
where we also bounded $\eta_t^{2}\le \eta_{\max}\eta_t$ and introduced $D:=\sup_t\norm{w_t-w^{\star}}$ to absorb the term arising from $\omega$ (see detailed derivation in the appendix). Equation (18) is exactly the bound (3) stated in Theorem~\ref{thm:main}.

\noindent\textbf{Step 14.} For a constant stepsize $\eta_t=\eta$, (18) reduces to (4). Choosing $\eta=\Theta(1/\sqrt{T})$ balances the two terms, yielding the rate $\mathcal{O}(1/\sqrt{T})$.

Thus the theorem is proved. \hfill $\square$

\section*{Rate Derivation (With Constants)}
From (4) with $\eta=\frac{c}{\sqrt{T}}$, $c\le\frac{1}{L(1+\omega)}$,
\[
\EE\!\big[F(\bar w_T)-F(w^{\star})\big]
\le \frac{\norm{w_0-w^{\star}}^{2}}{2c}\frac{1}{\sqrt{T}}
\;+\;
\frac{c}{2\sqrt{T}}(\sigma^{2}+\omega L^{2}D^{2})
\;=\;
\mathcal{O}\!\Big(\frac{1}{\sqrt{T}}\Big).
\]

\section*{Special Cases}
\begin{itemize}
\item \textbf{Exact gradients, no quantization} ($\sigma=0$, $\omega=0$): the bound collapses to $\frac{\norm{w_0-w^{\star}}^{2}}{2\eta T}$, recovering the deterministic gradient descent rate $\mathcal{O}(1/T)$.
\item \textbf{Only quantization noise} ($\sigma=0$, $\omega>0$): the variance term becomes $\frac{\eta\omega L^{2}D^{2}}{2}$, leading again to $\mathcal{O}(1/\sqrt{T})$ with optimal $\eta$.
\item \textbf{High‑precision communication} ($\omega\to0$): the theorem reduces to the classical SGD bound, confirming that QSGD is a strict generalisation.
\end{itemize}

\end{document}